\documentclass[11pt,a4paper]{report}
% Deutsch
\usepackage[german]{babel} % deutsch und deutsche Rechtschreibung
\usepackage[utf8]{inputenc} % Unicode Text 
\usepackage[T1]{fontenc} % Umlaute und deutsches Trennen
\usepackage{textcomp} % Euro
\usepackage[hyphens]{url}
% statt immer Ab\-schluss\-ar\-beit zu schreiben
% einfach hier sammeln mit -. 
\hyphenation{Ab-schluss-ar-beit}
% Vorsicht bei Umlauten und Bindestrichen
\hyphenation{Ver-st\"ar-ker-aus-gang}
 % eigene Hyphenations, die für das Dokument gelten
\usepackage{amssymb} % Symbole
\usepackage{emptypage} % Wirklich leer bei leeren Seiten
\usepackage{xcolor}

%% Fonts, ein kompletter Satz an Optionen

% Times New Roman, gewohnter Font, ok tt und serifenlos
%\usepackage{mathptmx} 
%\usepackage[scaled=.95]{helvet}
%\usepackage{courier}

% Palatino mit guten Fonts für tt und serifenlos
\usepackage{mathpazo} % Palatino, mal was anderes
\usepackage[scaled=.95]{helvet}
\usepackage{courier}

% New Century Schoolbook sieht auch nett aus (macht auch tt und serifenlos)
%\usepackage{newcent}

% Oder default serifenlos mit Helvetica 
% ich kann es nicht mehr sehen ...
%\renewcommand{\familydefault}{\sfdefault}

\usepackage{microtype}

% Bilder und Listings
\usepackage{graphicx} % wir wollen Bilder einfügen
\usepackage{subfig} % Teilbilder
\usepackage{wrapfig} % vielleicht doch besser vermeiden
\usepackage{listings} % schöne Quellcode-Listings
% ein paar Einstellungen für akzeptable Listings
\lstset{basicstyle=\ttfamily, columns=[l]flexible, mathescape=true, showstringspaces=false, numbers=left, numberstyle=\tiny}
\lstset{language=python} % und nur schöne Programmiersprachen ;-)
% und eine eigene Umgebung für Listings
\usepackage{float}
\newfloat{listing}{htbp}{scl}[chapter]
\floatname{listing}{Listing}

% Seitenlayout
\usepackage[paper=a4paper,width=14cm,left=35mm,height=22cm]{geometry}
\usepackage{setspace}
\linespread{1.15}
\setlength{\parskip}{0.5em}
\setlength{\parindent}{0em} % im Deutschen Einrückung nicht üblich, leider

% Seitenmarkierungen 
\newcommand{\phv}{\fontfamily{phv}\fontseries{m}\fontsize{9}{11}\selectfont}
\usepackage{fancyhdr} % Schickere Header und Footer
\pagestyle{fancy}
\renewcommand{\chaptermark}[1]{\markboth{#1}{}}
%\fancyhead[L]{\phv \leftmark}
\fancyhead[RE,LO]{\phv \nouppercase{\leftmark}}
\fancyhead[LE,RO]{\phv \thepage}
% Unten besser auf alles Verzichten
%\fancyfoot[L]{\textsf{\small \kurztitel}}
\fancyfoot[C]{\ } % keine Seitenzahl unten
%\fancyfoot[R]{\textsf{\small Medieninformatik}}

% Theorem-Umgebungen
\newtheorem{definition}{Definition}[chapter]
\newtheorem{satz}{Satz}[chapter]
\newtheorem{lemma}[satz]{Lemma} % gleicher Zähler wie Satz
\newtheorem{theorem}{Theorem}[chapter]
\newenvironment{beweis}[1][Beweis]{\begin{trivlist}
\item[\hskip \labelsep {\textit{#1 }}]}{\end{trivlist}}
\newcommand{\qed}{\hfill \ensuremath{\square}}

% Quellen teilen
\usepackage{bibtopic} 

% Spezialpakete
\usepackage{epigraph}
\setlength{\epigraphrule}{0pt} % kein Trennstrich

% damit wir nicht so viel tippen müssen, nur für Demo 
\usepackage{blindtext} 

% Zum Zeigen von Fehlern
\usepackage{soul}
\newcommand*\falsch{\st}


\usepackage{xr}
\externaldocument[T-]{thesis}
\begin{document}
\begin{center}
{\Large\bfseries Digitaler Anhang A}\\[0.5em]
{\large Einzelbegründungen der Entscheidungsmatrix}\\[1em]
zur Bachelorarbeit \glqq LLM-gestützte Rekonstruktion von
Softwareanforderungen aus der Entwicklungshistorie eines
Open-Source-Projekts\grqq{}\\
Omar Assad Salem, Hochschule RheinMain, 2026
\end{center}

\noindent Verweise auf Abschnitte, Tabellen und Abbildungen ohne den
Präfix A beziehen sich auf die Arbeit. Dieser Anhang war in der Fassung
vor der Kürzung Anhang~A der Arbeit; sein Inhalt ist unverändert. Ergänzt
sind am Ende die beiden Tabellen zum Quality-User-Story-Rahmenwerk aus
dem früheren Abschnitt~2.2.3, auf die die Begründungen verweisen.

\bigskip
\appendix

\chapter{Einzelbegründungen der Entscheidungsmatrix}
\label{cha:anhang-einzelbegruendungen}

Dieser Anhang gibt die Einzelbegründungen zu den Noten der
Entscheidungsmatrix aus Tabelle~\ref{T-tab:entscheidungsmatrix} in
gekürzter Form wieder. Die vollständigen Begründungen einschließlich
ihrer Änderungshistorie sind in der Arbeitsunterlage
\path{docs/formatauswahl.md}, Abschnitt~10, des mit der Arbeit
abgegebenen Repositorys dokumentiert; die Noten beider Fassungen
stimmen überein. Die Leitfragen der Kriterien sind in
Tabelle~\ref{T-tab:kriterienkatalog} aufgeführt.

Zu jeder Zelle werden drei Angaben gemacht. Die Belegstelle bezeichnet
den Quelleninhalt, auf den sich die Note stützt. Begründung und
Gegenargument sind eigene Einschätzungen; das Gegenargument
nennt den stärksten Einwand gegen die vergebene Note und die Note, zu
der er führen würde. Stützt sich ein Einwand auf Restprobleme, die eine
Quelle aus einer Erprobung berichtet, lässt er die Note nach der in
Abschnitt~\ref{T-subsec:auswahlverfahren} dargestellten Auslegung der
Stufe~2 unberührt; er wird dennoch aufgeführt. Der Wert \emph{nicht
anwendbar} ist keine schlechte Note, sondern zeigt an, dass eine Ebene
die betreffende Eigenschaft ihrem Zweck nach nicht regelt. Die Noten für
K6 und K9 beruhen auf der in Abschnitt~\ref{T-subsec:formulierungsprobe}
beschriebenen Messung, alle übrigen auf der Literatur.
Die Noten wurden nach Prüfung an den Quellen von Hand vergeben. Zur
Absicherung wurden Belegstellen, Notenvorschläge und Gegenargumente
eines Sprachmodells hinzugezogen.


\section{K1 Eindeutigkeit}
\label{sec:anhang-k1}

\begin{description}
  \item[(a) User Story, Note~1.]
  \emph{Beleg:} \cite[Abschn.~2, 2.3 und 3.1.6]{lucassen2016qus}.
  \emph{Begründung:} Die Schablone trennt Rolle, Means und Ends und legt
  für das Means drei grammatische Bestandteile fest, sodass zumindest
  ein Teil des Satzes einer festen Lesart zugeordnet ist. Der übrige
  Text bleibt frei formulierbar, und Eindeutigkeit wird als gesondertes
  Qualitätskriterium geführt.
  \emph{Gegenargument (spräche für Note~0):} Die Schablone ordnet den
  Satz, benennt aber keine Formulierung, die sie ausschließt.

  \item[(b) EARS, Note~2.]
  \emph{Beleg:} \cite[Abschn.~2, 4.1, 5 und 6]{mavin2009ears}.
  \emph{Begründung:} EARS legt die Bestandteile eines Satzes und ihre
  Reihenfolge fest und bestimmt, unter welcher Bedingung die
  Systemreaktion gilt; damit ist die Lesart des ganzen Satzes geregelt
  und nicht nur die eines Satzglieds. Die Quelle führt Mehrdeutigkeit
  als ausdrücklich adressiertes Problem und zeigt an Beispielen, welche
  Formulierungen ersetzt werden.
  \emph{Gegenargument (spräche für Note~1):} Nach Angabe der Autoren
  wurde Mehrdeutigkeit in ihrer Erprobung nur verringert, nicht
  beseitigt \cite[Abschn.~6]{mavin2009ears}.

  \item[(c) Normative Modalität, Note~1.]
  \emph{Beleg:} \cite[Abschn.~1--5]{rfc2119}; \cite[Abschn.~2]{rfc8174}.
  \emph{Begründung:} Für das Modalwort ist die Lesart durch ein
  geschlossenes Vokabular und die Regel der durchgängigen
  Großschreibung eindeutig festgelegt. Für den übrigen Satz trifft die
  Konvention keine Aussage.
  \emph{Gegenargument (spräche für Note~2):} Innerhalb ihrer Reichweite
  ist die Regel vollständig und an keine Bedingung geknüpft. Als nicht
  anwendbar wird das Kriterium nicht geführt, weil die Ebene Lesarten
  regelt, wenn auch nur für ein einzelnes Wort.

  \item[(d) Qualitäts-, Sprach- und Attributregeln, Note~2.]
  \emph{Beleg:} \cite[Abschn.~5.2.5 und 5.2.7]{iso29148};
  \cite[Abschn.~Testing the Hypothesis]{mavin2010bigears};
  \cite[Tab.~1 und Abschn.~3.1.6]{lucassen2016qus}.
  \emph{Begründung:} Abschnitt~5.2.7 des Standards führt in neun
  Aufzählungspunkten Wortklassen auf, die in Anforderungen zu vermeiden
  sind. Die Ebene schließt damit als einzige der vier Ebenen einzelne
  Formulierungen namentlich und in normativer Form aus.
  \emph{Gegenargument (spräche für Note~1):} Die Regeln benennen
  Wortklassen, aber kein Verfahren, das Eindeutigkeit herstellt; sie
  machen Verstöße erkennbar, verhindern sie aber nicht.
\end{description}


\section{K2 Formale Struktur}
\label{sec:anhang-k2}

Für die Ebenen~(c) und~(d) wird das Kriterium als nicht anwendbar
geführt und nicht mit~0 bewertet, weil beide ihrem Zweck nach keine
Satzbestandteile regeln. Dieselbe Erwägung liegt der Nichtanwendbarkeit
von K6 für diese Ebenen zugrunde. Dagegen spricht, dass ISO/IEC/IEEE~29148
eine Schlüsselwortkonvention mit einem Aufbausatz verbindet
\cite[Abschn.~5.2.4]{iso29148} und dass die Zuordnung dieses Aufbausatzes
zu Ebene~(b) eine eigene Setzung ist. Die Einstufung als nicht
anwendbar ist insoweit auch eine Folge dieser Zuordnung.

\begin{description}
  \item[(a) User Story, Note~1.]
  \emph{Beleg:} \cite[Abschn.~2 mit Abb.~1, 2.1, 2.3, 3.2.2 und
  4.7]{lucassen2016qus}.
  \emph{Begründung:} Das Begriffsmodell benennt vier Bestandteile und
  legt für das Means drei grammatische Elemente fest, sodass ein Satz in
  der verbreiteten Schablone ohne Deutung zerlegbar ist. Die Schablone
  selbst gibt der Ansatz jedoch nicht vor; er verweist auf über achtzig
  syntaktische Varianten und bestimmt die maßgebliche Form je Datensatz
  nach ihrer Häufigkeit.
  \emph{Gegenargument (spräche für Note~2):} Die Kriterien Well-formed,
  Atomic und Minimal (Tabelle~\ref{tab:qus-kriterien}) legen fest, welche
  Bestandteile ein Satz enthalten muss und welche nicht; in der
  verbreiteten Form steht auch die Reihenfolge fest.
  \emph{Hinweis zur Stufe~2:} Die geringe Genauigkeit der selbsttätigen
  Prüfung beim Kriterium Well-formed (Tabelle~\ref{tab:aqusa}) ist ein
  Erprobungsbefund und trägt die Note nicht. Diese stützt sich allein
  darauf, dass der Ansatz die Schablone offenlässt.

  \item[(b) EARS, Note~2.]
  \emph{Beleg:} \cite[Abschn.~4.1--4.7]{mavin2009ears};
  \cite[Abschn.~EARS Templates]{mavin2010bigears}.
  \emph{Begründung:} EARS benennt vier Bestandteile und erklärt ihre
  Reihenfolge ausdrücklich für bedeutsam, weil sie der zeitlichen Logik
  folgt. Vier der fünf Muster geben zusätzlich eine eigene Satzform mit
  festem Schlüsselwort vor; das fünfte gilt ohne Bedingung und ohne
  Auslöser.
  \emph{Gegenargument (spräche für Note~1):} Vorbedingung und Auslöser
  sind optional, und die Abgrenzung zwischen den Mustern Ubiquitous und
  Event-driven kann nach Angabe der Autoren der Folgearbeit schwierig
  sein \cite[Abschn.~Lessons Learned Applying EARS]{mavin2010bigears}.
  \emph{Hinweis zur Stufe~2:} Der zweite Teil des Einwands ist ein
  Erprobungsbefund und trägt die Note nicht. Der erste Teil betrifft die
  Reichweite der Regel; optionale Satzteile heben die festgelegte
  Reihenfolge jedoch nicht auf, sie entfallen nur.

  \item[(c) Normative Modalität, nicht anwendbar.]
  \emph{Beleg:} \cite[Abschn.~1--5]{rfc2119} und
  \cite[Abschn.~2]{rfc8174} treffen keine Aussage zum Satzbau.
  \emph{Begründung:} Die Ebene legt allein die Bedeutung von
  Schlüsselwörtern fest und regelt weder Satzbestandteile noch deren
  Reihenfolge.
  \emph{Gegenargument (spräche für Note~0):} siehe die Vorbemerkung zu
  diesem Kriterium.

  \item[(d) Qualitäts-, Sprach- und Attributregeln, nicht anwendbar.]
  \emph{Beleg:} \cite[Abschn.~5.2.7, 5.2.8.1 und 5.2.8.2]{iso29148};
  \cite[Tab.~1]{lucassen2016qus};
  \cite[Abschn.~Testing the Hypothesis]{mavin2010bigears}.
  \emph{Begründung:} Die Ebene beurteilt eine bereits formulierte
  Anforderung. Abschnitt~5.2.7 nennt zu vermeidende Wortklassen,
  Abschnitt~5.2.8 Attribute neben dem Satz; beides sind weder
  Satzbestandteile noch eine Reihenfolge.
  \emph{Gegenargument (spräche für Note~0):} siehe die Vorbemerkung zu
  diesem Kriterium. Die Attribute sind zwar benannte Felder mit fester
  Bedeutung, stehen aber neben dem Satz und nicht in ihm; erheblich
  werden sie erst für K4.
\end{description}


\section{K3 Verifizierbarkeit}
\label{sec:anhang-k3}

\begin{description}
  \item[(a) User Story, Note~1.]
  \emph{Beleg:} \cite[Abschn.~1, 3.1.3 und Tab.~1]{lucassen2016qus}.
  \emph{Begründung:} Der Ansatz zählt Akzeptanzkriterien ausdrücklich zu
  den drei Bestandteilen einer User Story. Im Anforderungstext selbst
  dürfen sie nach dem Kriterium Minimal jedoch nicht stehen, und keines
  der dreizehn Qualitätskriterien prüft die Nachweisbarkeit.
  \emph{Gegenargument (spräche für Note~0):} Bewertet wird der
  Anforderungstext, und für ihn erzwingt der Ansatz keine überprüfbare
  Angabe; die Akzeptanzkriterien liegen außerhalb und bleiben
  ungeprüft.

  \item[(b) EARS, Note~1.]
  \emph{Beleg:} \cite[Abschn.~2 und 4.1]{mavin2009ears};
  \cite[Abschn.~Residual Problems]{mavin2010bigears}.
  \emph{Begründung:} Die Syntax erzwingt eine Systemreaktion und damit
  ein beobachtbares Verhalten, und Untestbarkeit gehört zu den acht
  Problemen, die EARS ausdrücklich adressiert. Eine Regel, die eine
  messbare Größe oder ein Abnahmekriterium verlangt, enthält der Ansatz
  jedoch nicht.
  \emph{Gegenargument (spräche für Note~2):} Die Beseitigung der
  Untestbarkeit ist ein erklärtes Ziel des Ansatzes, und Vorbedingung,
  Auslöser und Systemreaktion benennen genau die Angaben, gegen die
  geprüft wird.
  \emph{Hinweis zur Stufe~2:} Die in der Folgearbeit berichtete
  verbleibende Untestbarkeit durch nicht quantifizierbare Wörter ist ein
  Erprobungsbefund und trägt die Note nicht. Diese stützt sich allein
  darauf, dass eine Regel zur Quantifizierung fehlt.

  \item[(c) Normative Modalität, nicht anwendbar.]
  \emph{Beleg:} \cite{rfc2119} und \cite{rfc8174} treffen keine Aussage
  darüber, woran die Erfüllung einer Anforderung überprüft wird.
  \emph{Begründung:} Die Konvention legt allein die
  Verbindlichkeitsstufe fest und regelt ihrem Zweck nach nicht die
  Nachweisbarkeit.
  \emph{Gegenargument (spräche für Note~0):} Die Verbindlichkeitsstufe
  entscheidet mit darüber, ob etwas überhaupt nachzuweisen ist; die
  Konvention berührt die Nachweisbarkeit damit mittelbar.

  \item[(d) Qualitäts-, Sprach- und Attributregeln, Note~2.]
  \emph{Beleg:} \cite[Abschn.~5.2.5 und 5.2.7]{iso29148}.
  \emph{Begründung:} Der Standard führt die Nachweisbarkeit als eigenes
  Merkmal jeder einzelnen Anforderung und begründet die Liste zu
  vermeidender Ausdrücke ausdrücklich damit, dass diese die
  Überprüfung erschweren oder unmöglich machen. Damit liegt eine
  ausdrückliche und weitgehend operationalisierte Regel vor.
  \emph{Gegenargument (spräche für Note~1):} Messbarkeit wird nur als
  förderlich für die Nachweisbarkeit bezeichnet, nicht verlangt, und die
  Regeln nennen, was zu vermeiden ist, nicht, welche Angabe eine
  Anforderung enthalten muss.
\end{description}


\section{K4 Maschinelle Verarbeitbarkeit}
\label{sec:anhang-k4}

Das Kriterium ist eine eigene Setzung, gestützt auf die Klauselstruktur
von EARS (Tabelle~\ref{T-tab:kriterienkatalog}). Ein aus einem der
bewerteten Ansätze abgeleitetes Kriterium begünstigt diesen Ansatz; die
Note der Ebene~(b) ist daher nur begrenzt aussagekräftig. Dieser
Vorbehalt wurde vor der Notenvergabe festgehalten und wird in
Kapitel~\ref{T-cha:diskussion} als Bedrohung der Konstruktvalidität
aufgegriffen.

\begin{description}
  \item[(a) User Story, Note~2.]
  \emph{Beleg:} \cite[Abschn.~2 mit Abb.~1, 2.3, 3.2 und
  4.3]{lucassen2016qus}.
  \emph{Begründung:} Der Ansatz definiert eine User Story formal als
  Viertupel aus Rolle, Means, Ends und Format und das Means seinerseits
  als Fünftupel. Damit besteht eine ausdrückliche Zuordnung von Satzteil
  zu Feld bis hinunter zur grammatischen Ebene des Means.
  \emph{Gegenargument (spräche für Note~1):} Das Format ist selbst eines
  der Felder und wird nicht festgelegt. Ohne Kenntnis der verwendeten
  Schablone bestimmt der Ansatz allein nicht, wo Rolle, Means und Ends im
  Text beginnen.

  \item[(b) EARS, Note~2.]
  \emph{Beleg:} \cite[Abschn.~4.1--4.6]{mavin2009ears};
  \cite[Tab.~I]{wilkinson2018mitigation}.
  \emph{Begründung:} Die generische Syntax besteht aus benannten
  Platzhaltern, die im Satz durch feste Schlüsselwörter voneinander
  abgegrenzt sind. Eine spätere Arbeit derselben Autoren behandelt diese
  Bestandteile bereits als einzeln adressierbare Teile.
  \emph{Gegenargument (spräche für Note~1):} Die Quellen benennen
  Satzteile, definieren aber weder ein Datenmodell noch Feldnamen; die
  Abbildung auf Datenfelder ist eine Ergänzung der vorliegenden Arbeit.

  \item[(c) Normative Modalität, Note~1.]
  \emph{Beleg:} \cite[Abschn.~1--5]{rfc2119}; \cite[Abschn.~2]{rfc8174}.
  \emph{Begründung:} Das Schlüsselwort lässt sich verlustfrei auf ein
  einzelnes Feld mit fünf zulässigen Stufen abbilden, und die Regel der
  Großschreibung grenzt es im Text eindeutig ab. Abgedeckt ist damit
  genau ein Feld, nicht der Satz.
  \emph{Gegenargument (spräche für nicht anwendbar):} Die Ebene bringt
  keinen Satz hervor, während die Leitfrage von einem nach dem Ansatz
  formulierten Satz spricht. Als nicht anwendbar wird das Kriterium
  dennoch nicht geführt, weil die Ebene etwas Abbildbares regelt, wenn
  auch eng.

  \item[(d) Qualitäts-, Sprach- und Attributregeln, Note~1.]
  \emph{Beleg:} \cite[Abschn.~5.2.8.1 und 5.2.8.2]{iso29148}.
  \emph{Begründung:} Der Standard benennt Attribute wie Kennung,
  Versionsnummer, Eigner, Priorität, Risiko, Begründung und
  Schwierigkeitsgrad als Felder mit fester Bedeutung, die eine
  Anforderung begleiten. Sie stehen jedoch neben dem Satz, und eine
  Zuordnung von Satzteil zu Feld gibt der Ansatz nicht vor.
  \emph{Gegenargument (spräche für Note~2):} Für die Ausgabe der
  Pipeline sind gerade diese Attribute die Felder neben dem
  Anforderungstext; ihre Benennung ist im Standard festgelegt und
  vollständig übernehmbar.
\end{description}


\section{K5 Eignung für Modellausgaben}
\label{sec:anhang-k5}

Die Noten dieses Kriteriums stützen sich nicht auf Untersuchungen zum
Verhalten von Sprachmodellen, sondern auf drei strukturelle
Ersatzgrößen: die Zahl der Pflichtbestandteile, die Zahl zulässiger
Schreibvarianten und das Vorhandensein eines Musters für den Fall, dass
keine Anforderung vorliegt. Sie sind vorläufig und werden nach der
Entwicklung der Prompts überprüft.

\begin{description}
  \item[(a) User Story, Note~1.]
  \emph{Beleg:} \cite[Tab.~1, Abschn.~2 und 3.2.2]{lucassen2016qus}.
  \emph{Begründung:} Zwei Pflichtbestandteile und die Begrenzung durch
  das Kriterium Minimal sind ein wirksamer Mechanismus zur Steuerung der
  Ausgabe. Der Ansatz legt jedoch keine Schablone fest, lässt über
  achtzig Varianten zu und kennt kein Muster für den Fall, dass keine
  Anforderung vorliegt.
  \emph{Gegenargument (spräche für Note~0):} Die Ebene ist nur in vier
  von zwölf gemessenen Fällen gedeckt, ausschließlich wegen der fehlenden Rolle (Tabelle~\ref{T-tab:deckung}). Ein Format, das einen unbelegten
  Bestandteil verlangt, drängt ein Modell zum Erfinden.

  \item[(b) EARS, Note~2.]
  \emph{Beleg:} \cite[Abschn.~4.1--4.6]{mavin2009ears};
  \cite[Abschn.~Method]{mavin2010bigears}.
  \emph{Begründung:} Zwei Pflichtbestandteile und ein geschlossener Satz
  von fünf Mustern mit festgelegten Schlüsselwörtern halten die Zahl
  zulässiger Schreibvarianten sehr klein. Für den Fall, dass keine
  Anforderung vorliegt, beschreibt die Folgearbeit ein ausdrückliches
  Verfahren: Jede Aussage wird vor der Umformulierung als Anforderung,
  Überschrift, Information oder für das betrachtete System nicht
  anwendbar eingestuft.
  \emph{Gegenargument (spräche für Note~1):} Diese Einstufung gehört zum
  Verfahren der Folgearbeit und nicht zur Syntax; EARS als Format kennt
  kein Muster für eine Aussage, die keine Anforderung ist.

  \item[(c) Normative Modalität, Note~1.]
  \emph{Beleg:} \cite[Abschn.~1--6]{rfc2119}; \cite[Abschn.~2]{rfc8174}.
  \emph{Begründung:} Ein geschlossenes Vokabular aus fünf Stufen mit
  eindeutiger Schreibweise ist für ein Sprachmodell leicht einzuhalten.
  Die in Abschnitt~6 von RFC~2119 geforderte sparsame Verwendung
  verlangt jedoch eine Abwägung, die ein Modell aus den Quelldaten nicht
  treffen kann.
  \emph{Gegenargument (spräche für Note~0):} Eine Verbindlichkeitsstufe
  ist in keinem der zwölf gemessenen Fälle belegt
  (Tabelle~\ref{T-tab:deckung}); jede rekonstruierte Anforderung erforderte
  damit eine erfundene Stufe.

  \item[(d) Qualitäts-, Sprach- und Attributregeln, nicht anwendbar.]
  \emph{Beleg:} \cite[Abschn.~5.2.5--5.2.8]{iso29148};
  \cite[Tab.~1]{lucassen2016qus};
  \cite[Abschn.~Testing the Hypothesis]{mavin2010bigears}.
  \emph{Begründung:} Die Ebene gibt weder ein Ausgabemuster noch
  Pflichtbestandteile vor, während alle drei Ersatzgrößen ein solches
  Muster voraussetzen.
  \emph{Gegenargument (spräche für eine Note):} Die Liste zu vermeidender
  Ausdrücke aus Abschnitt~5.2.7 lässt sich unmittelbar in eine Anweisung
  an ein Sprachmodell übersetzen. Für dessen Steuerung ist die Ebene
  damit einschlägig, nur nicht über ein Satzmuster.
\end{description}


\section{K6 Eignung für die drei Anforderungsarten}
\label{sec:anhang-k6}

Die Noten dieses Kriteriums stammen nicht aus der Literatur, sondern aus
der in Abschnitt~\ref{T-subsec:formulierungsprobe} beschriebenen Messung.
Rahmen und Referenzinhalte wurden von Hand festgelegt, die sechs
Sätze der Tabelle~\ref{tab:anhang-k6-versuche} von Hand formuliert
und eingestuft. Zur Absicherung hat ein Sprachmodell
Vorschläge für Rahmen und Referenzinhalte geliefert und die
Einstufungen anschließend gegengeprüft; eine Note hat es nicht
vorgeschlagen. Die vollständige Dokumentation der Messung
steht in \path{docs/formatprobe_protokoll.md}.

Für die Ebenen~(c) und~(d) wird das Kriterium als nicht anwendbar
geführt, weil ein Formulierungsversuch dort den Satzbau einer anderen
Ebene übernehmen müsste. Dagegen spricht, dass die beiden gemessenen
Kriterien damit nur dort greifen, wo eine Ebene ein Satzmuster vorgibt;
Ebene~(d) hat deshalb weder bei K6 noch bei K9 eine gemessene Note.
Beide bewerteten Ebenen beziehen zudem je einen Pflichtbestandteil aus
dem einheitlichen Rahmen der Probe, Ebene~(a) die Rolle und Ebene~(b)
den Systemnamen; das Kriterium unterscheidet die beiden Ebenen daher
nicht.

\begin{table}[htbp]
  \centering
  \caption{Formulierungsversuche der Probe für das Kriterium~K6}
  \label{tab:anhang-k6-versuche}
  \begin{tabular}{@{}>{\raggedright\arraybackslash}p{2.2cm}%
                    >{\raggedright\arraybackslash}p{2.2cm}%
                    >{\raggedright\arraybackslash}p{7.4cm}@{}}
    \toprule
    \textbf{Art} & \textbf{Ebene} & \textbf{Anforderungssatz} \\
    \midrule
    funktional & (a) User Story &
    As a developer, I want inputs to be checked for validity before they
    are processed. \\
    funktional & (b) EARS &
    The library shall check inputs for validity before processing
    them. \\
    \addlinespace
    leistungs\-bezogen & (a) User Story &
    As a developer, I want 95\,\% of transactions to be processed in
    less than 1\,s. \\
    leistungs\-bezogen & (b) EARS &
    The library shall process 95\,\% of transactions in less than
    1\,s. \\
    \addlinespace
    sicherheits\-bezogen & (a) User Story &
    As a developer, I want the integrity of data in critical variables
    to be checked. \\
    sicherheits\-bezogen & (b) EARS &
    The library shall check the integrity of data in critical
    variables. \\
    \bottomrule
  \end{tabular}
\end{table}

\begin{description}
  \item[(a) User Story, Note~2.]
  \emph{Beleg:} drei von drei Versuchen gelungen
  (Tabelle~\ref{tab:anhang-k6-versuche}).
  \emph{Begründung:} In allen drei Versuchen sind Rolle und Means
  besetzt, ohne dass eine Angabe hinzukommt, die Rahmen oder
  Referenzinhalt nicht hergeben. Eine von der Anforderungsart abhängige
  Schwäche zeigt sich bei keiner der drei Arten.
  \emph{Gegenargument (spräche für Note~0):} Die Rolle stammt in allen
  drei Sätzen aus dem Rahmen und nicht aus dem Referenzinhalt; würde sie
  nicht als belegt gelten, gelänge kein einziger Versuch.

  \item[(b) EARS, Note~2.]
  \emph{Beleg:} drei von drei Versuchen gelungen
  (Tabelle~\ref{tab:anhang-k6-versuche}).
  \emph{Begründung:} In allen drei Versuchen sind Systemname und
  Systemreaktion besetzt, wobei die Systemreaktion jeweils dem
  Referenzinhalt selbst entstammt. Auch hier zeigt sich bei keiner der
  drei Anforderungsarten eine Schwäche.
  \emph{Gegenargument (spräche für Note~0):} Der Systemname stammt
  ebenso aus dem Rahmen wie bei Ebene~(a) die Rolle; würde auch er nicht
  als belegt gelten, gelänge gleichfalls kein Versuch.

  \item[(c) Normative Modalität, nicht anwendbar.]
  \emph{Beleg:} \cite[Abschn.~1--5]{rfc2119} und
  \cite[Abschn.~2]{rfc8174} treffen keine Aussage zum Satzbau.
  \emph{Begründung:} Die Ebene legt allein die Bedeutung von
  Schlüsselwörtern fest; ein Formulierungsversuch müsste den Satzbau
  einer anderen Ebene übernehmen und bewertete damit diese.
  \emph{Gegenargument (spräche für eine Note):} siehe die Vorbemerkung
  zu diesem Kriterium.

  \item[(d) Qualitäts-, Sprach- und Attributregeln, nicht anwendbar.]
  \emph{Beleg:} \cite[Abschn.~5.2.5--5.2.7]{iso29148}.
  \emph{Begründung:} Die Ebene nennt Merkmale und sprachliche Kriterien,
  aber keine Satzbestandteile, und beurteilt eine bereits formulierte
  Anforderung.
  \emph{Gegenargument (spräche für eine Note):} Ob die Qualitätsregeln
  bei den drei Anforderungsarten unterschiedlich greifen, bleibt damit
  ungemessen; dieser Gesichtspunkt geht nur mittelbar über K3 ein.
\end{description}


\section{K7 Verständlichkeit}
\label{sec:anhang-k7}

Dieses Kriterium wird nur ersatzweise bewertet. Belastbar wird es erst
durch eine Bewertung durch das Qrisp-Team, die für die Zeit nach der
Abgabe vorgesehen ist; die Selbstbewertung in
Abschnitt~\ref{T-sec:eval-bewertung} ersetzt sie nicht. Sein zweiter
Prüfpunkt, der Abstand zur Ausdrucksweise in den Pull Requests, ist eine
eigene Setzung ohne Beleg in einer Quelle; wo er eine Note
trägt, ist das bei der betreffenden Zelle ausgewiesen.

\begin{description}
  \item[(a) User Story, Note~1.]
  \emph{Beleg:} \cite[Abschn.~1 und 2.2]{lucassen2016qus}.
  \emph{Begründung:} Drei Signalphrasen in Alltagsenglisch ohne
  Fachvokabular verlangen einer lesenden Person kaum Vorwissen ab. Die
  Schablone setzt jedoch eine Rolle aus der Anwendungsdomäne voraus,
  während die Pull Requests des Projekts sich an andere Entwickelnde
  richten und keinen Nutznießer benennen.
  \emph{Gegenargument (spräche für Note~2):} Nach dem ersten Prüfpunkt,
  der Zahl der zu lernenden Schlüsselwörter und Regeln, schneidet die
  Ebene von allen vieren am besten ab.
  \emph{Ausweis der eigenen Setzung:} Der Abzug beruht allein auf dem
  zweiten Prüfpunkt und ist damit gesetzt, nicht belegt.

  \item[(b) EARS, Note~2.]
  \emph{Beleg:} \cite[Abschn.~2 und 4.1--4.6]{mavin2009ears}.
  \emph{Begründung:} Der Ansatz lehnt formale und grafische Notationen
  ausdrücklich wegen Übersetzungsfehlern, Sprachbarriere und
  Schulungsaufwand ab und ist darauf angelegt, ohne Schulung lesbar zu
  bleiben. Seine Satzform nennt das System als Handelnden, so wie die
  Pull Requests über die Bibliothek sprechen.
  \emph{Gegenargument (spräche für Note~1):} Zu lernen sind fünf Muster
  mit vier Schlüsselwörtern, das Wort \emph{shall} und die verbindliche
  Klauselreihenfolge; \emph{shall} ist in der Sprache der Pull Requests
  unüblich.
  \emph{Hinweis zur Stufe~2:} Die in der Folgearbeit berichtete
  Schwierigkeit, die Muster Ubiquitous und Event-driven auseinander zu
  halten, ist ein Erprobungsbefund und trägt die Note nicht
  \cite[Abschn.~Lessons Learned Applying EARS]{mavin2010bigears}.

  \item[(c) Normative Modalität, Note~1.]
  \emph{Beleg:} \cite[Abschn.~1--5]{rfc2119};
  \cite[Abschn.~1 und 2]{rfc8174}.
  \emph{Begründung:} Die fünf Verbindlichkeitsstufen werden mit elf
  geläufigen englischen Wörtern ausgedrückt, sodass dem Wortschatz nach
  nichts Fremdes zu lernen ist. Genau darin liegt die Falle: Wer die
  Konvention nicht kennt, liest \emph{SHOULD} als unverbindliche
  Empfehlung, und RFC~8174 wurde eigens erlassen, weil die
  Unterscheidung nach Groß- und Kleinschreibung nicht eindeutig war.
  \emph{Gegenargument (spräche für Note~2):} Fünf Bedeutungen sind ein
  überschaubarer Lernaufwand, und die durchgängige Großschreibung macht
  sichtbar, dass überhaupt eine Konvention gilt.

  \item[(d) Qualitäts-, Sprach- und Attributregeln, Note~2.]
  \emph{Beleg:} \cite[Abschn.~5.2.6 und 5.2.7]{iso29148}.
  \emph{Begründung:} Der Standard führt die Verständlichkeit als eigenes
  Merkmal und untersagt unbestimmte und mehrdeutige Ausdrücke. Die
  Regeln richten sich an die verfassende Person und nicht an die
  lesende; diese muss nichts lernen.
  \emph{Gegenargument (spräche für Note~1):} Das Merkmal gilt nach
  Abschnitt~5.2.6 für die Menge der Anforderungen, während das Kriterium
  nach dem Lesen eines einzelnen Satzes fragt; überdies sagt der
  Standard, was zu vermeiden ist, und nicht, wie ein verständlicher Satz
  aussieht.
\end{description}


\section{K8 Automatische Validierbarkeit}
\label{sec:anhang-k8}

\begin{description}
  \item[(a) User Story, Note~1.]
  \emph{Beleg:} \cite[Abschn.~4, 4.2 und 4.7]{lucassen2016qus}.
  \emph{Begründung:} Mit AQUSA liegt ein veröffentlichtes Werkzeug vor,
  das die Formatkonformität je Kriterium prüft; ein Mechanismus ist
  damit zweifelsfrei vorhanden. Er ist jedoch an eine Bedingung
  geknüpft, weil der Ansatz keine feste Schablone vorgibt und deshalb
  gegen eine je Datensatz aus der Häufigkeit bestimmte Form geprüft
  wird.
  \emph{Gegenargument (spräche für Note~2):} Dies ist die einzige der
  vier Ebenen mit einem tatsächlich gebauten und veröffentlichten
  Prüfwerkzeug.
  \emph{Hinweis zur Stufe~2:} Die gemessene Leistungsfähigkeit dieses
  Werkzeugs (Tabelle~\ref{tab:aqusa}) ist ein Erprobungsbefund und trägt
  die Note nicht.

  \item[(b) EARS, Note~2.]
  \emph{Beleg:} \cite[Abschn.~4.1--4.6]{mavin2009ears}.
  \emph{Begründung:} Die Schlüsselwörter der vier bedingten Muster und
  das Modalwort \emph{shall} bilden ein geschlossenes Vokabular, das ein
  Programm unmittelbar suchen kann, und das fünfte Muster ist am Fehlen
  eines Schlüsselworts erkennbar. Die festgelegte Klauselreihenfolge
  macht zusätzlich prüfbar, ob ein Satz dem Muster folgt und nicht nur
  dessen Wörter enthält.
  \emph{Gegenargument (spräche für Note~1):} In den ausgewerteten
  Quellen ist kein Prüfwerkzeug für den Ansatz beschrieben; die
  maschinelle Prüfbarkeit ist aus der Syntax abgeleitet, aber nicht
  gezeigt.

  \item[(c) Normative Modalität, Note~2.]
  \emph{Beleg:} \cite[Abschn.~1--6]{rfc2119}; \cite[Abschn.~2]{rfc8174}.
  \emph{Begründung:} Das Vokabular ist geschlossen und durch die Regel
  der durchgängigen Großschreibung im Fließtext eindeutig abgegrenzt,
  sodass eine einfache Suche genügt, um festzustellen, ob und welche
  Stufe verwendet wurde.
  \emph{Gegenargument (spräche für Note~1):} Ein Programm findet das
  Schlüsselwort, kann aber die in Abschnitt~6 geforderte sparsame
  Verwendung nicht prüfen; diese Regel für inhaltlich zentral zu halten,
  ist allerdings eine eigene Einordnung, da die Quelle den
  Abschnitt als Hinweis zur Anwendung führt.

  \item[(d) Qualitäts-, Sprach- und Attributregeln, Note~1.]
  \emph{Beleg:} \cite[Abschn.~5.2.5 und 5.2.7]{iso29148}.
  \emph{Begründung:} Die neun Klassen unbestimmter Ausdrücke lassen sich
  in Suchregeln überführen, ein Mechanismus ist also vorhanden. Er
  reicht aber nur so weit: Die Wortbeispiele je Klasse sind nicht
  abschließend, und Merkmale wie notwendig, umsetzbar oder korrekt sind
  maschinell nicht prüfbar.
  \emph{Gegenargument (spräche für Note~2):} Für den Zweck des
  Kriteriums genügt, dass die sprachlichen Regeln prüfbar sind; die
  übrigen Merkmale betreffen den Inhalt und nicht die Formkonformität.
\end{description}


\section{K9 Deckung durch die Quelldaten}
\label{sec:anhang-k9}

Die Noten dieses Kriteriums stammen aus der in
Abschnitt~\ref{T-subsec:formulierungsprobe} beschriebenen Messung an zwölf
Pull Requests; Tabelle~\ref{T-tab:deckung} führt die gemessenen Werte auf.
Als gedeckt gilt ein Fall, wenn sich jeder Pflichtbestandteil der Ebene
aus Titel, Beschreibung, Kommentaren oder Änderungsdiff belegen lässt,
bei Ebene~(a) also Rolle und Means, bei Ebene~(b) Systemname und
Systemreaktion und bei Ebene~(c) die Verbindlichkeitsstufe. Bei zwölf
Fällen sind die Schwellen von zehn gedeckten Fällen für die Note~2 und
von fünf für die Note~1 wirksam.

Die Messung wurde von Hand durchgeführt; zur Absicherung hat ein
Sprachmodell alle zwölf Fälle unabhängig gegengeprüft. Die fünf
Einstufungen, die daraufhin nach Prüfung an der Quelle geändert
wurden, betrafen sämtlich Ebene~(c). Die vollständige
Dokumentation steht in \path{docs/formatprobe_protokoll.md}.

\begin{description}
  \item[(a) User Story, Note~0.]
  \emph{Beleg:} vier von zwölf Fällen gedeckt
  (Tabelle~\ref{T-tab:deckung}).
  \emph{Begründung:} Alle acht nicht gedeckten Fälle scheitern allein am
  Bestandteil Rolle, während das Means in allen zwölf Fällen belegt ist.
  Zwei der vier gedeckten Fälle nennen eine Personengruppe nur pauschal;
  ohne sie wären zwei von zwölf Fällen gedeckt, und die Note bliebe
  bei~0.
  \emph{Gegenargument (spräche für Note~1):} Zur nächsten Schwelle fehlt
  ein einziger Fall.

  \item[(b) EARS, Note~1.]
  \emph{Beleg:} neun von zwölf Fällen gedeckt
  (Tabelle~\ref{T-tab:deckung}).
  \emph{Begründung:} In den drei nicht gedeckten Fällen fehlt entweder
  ein eindeutig benanntes Bauteil oder eine konkrete Handlung. In zwei
  von ihnen lag der Änderungstext nur teilweise vor, und nicht sichtbare
  Einzelheiten wurden nicht ergänzt.
  \emph{Gegenargument (spräche für Note~2):} Zur nächsten Schwelle fehlt
  ein einziger Fall.

  \item[(c) Normative Modalität, Note~0.]
  \emph{Beleg:} kein Fall von zwölf gedeckt
  (Tabelle~\ref{T-tab:deckung}).
  \emph{Begründung:} Modalaussagen kommen in fünf Fällen vor, betreffen
  dort aber das Vorgehen im Pull Request, eine andere Anforderung oder
  ein benachbartes Bauteil. Als Beleg gilt eine Stufe nur, wenn sich die
  Aussage auf die rekonstruierte Anforderung selbst bezieht.
  \emph{Gegenargument (spräche für Note~1):} Ohne diese Einschränkung
  wären fünf von zwölf Fällen gedeckt und die Note läge bei~1, sodass
  auch die oben genannte Ausschlussregel nicht mehr griffe; die fünf
  Einstufungen gehen sämtlich auf die Gegenprüfung durch ein
  Sprachmodell zurück.

  \item[(d) Qualitäts-, Sprach- und Attributregeln, nicht anwendbar.]
  \emph{Beleg:} \cite[Abschn.~5.2.5--5.2.7]{iso29148}.
  \emph{Begründung:} Die Ebene schreibt keine Satzbestandteile vor,
  sondern Qualitäts- und Sprachregeln; ein Deckungsanteil ist für sie
  nicht definiert.
  \emph{Gegenargument (spräche für eine Note):} Der Standard sieht
  Attribute neben dem Anforderungssatz vor
  \cite[Abschn.~5.2.8.1]{iso29148}, und das gewählte Zielformat verlangt
  drei von ihnen; ob die Quelldaten diese hergeben, wurde nicht
  gemessen.
\end{description}


\section{K10 Kombinierbarkeit}
\label{sec:anhang-k10}

\begin{description}
  \item[(a) User Story, Note~1.]
  \emph{Beleg:} \cite[Abschn.~5.2.4 NOTE~2]{iso29148};
  \cite[Abschn.~1 und Tab.~1]{lucassen2016qus}.
  \emph{Begründung:} Der Standard lässt User Stories ausdrücklich zu,
  sodass die Verbindung mit Ebene~(d) belegt möglich ist. Die Schablone
  besetzt jedoch den ganzen Satz und lässt nach dem Kriterium Minimal
  keinen Platz für die Klauseln der Ebene~(b) oder ein Modalwort der
  Ebene~(c).
  \emph{Gegenargument (spräche für Note~2):} Die Ebene ist die einzige,
  deren Verträglichkeit eine Norm ausdrücklich bestätigt, und die
  Unvereinbarkeit mit Ebene~(b) betrifft nur den einzelnen Satz.

  \item[(b) EARS, Note~1.]
  \emph{Beleg:} \cite[Abschn.~4.1]{mavin2009ears};
  \cite[Abschn.~5.2.4]{iso29148}; \cite[Abschn.~1--5]{rfc2119}.
  \emph{Begründung:} Mit Ebene~(d) ist der Ansatz verträglich, weil er
  dasselbe Modalwort verwendet, das der Standard für verbindliche
  Anforderungen vorsieht. Zur Ebene~(c) in ihrer Ausprägung nach
  RFC~2119 besteht dagegen ein Widerspruch, weil der Ansatz nur
  \emph{shall} kennt und keine Abstufung.
  \emph{Gegenargument (spräche für Note~2):} Der Widerspruch besteht
  nicht zur Ebene~(c) als solcher, sondern nur zu einer ihrer beiden
  Ausprägungen; mit den Schlüsselwörtern des Standards ist der Ansatz
  widerspruchsfrei.

  \item[(c) Normative Modalität, Note~1.]
  \emph{Beleg:} \cite[Abschn.~1--5]{rfc2119};
  \cite[Abschn.~5.2.4]{iso29148}.
  \emph{Begründung:} Der Konflikt liegt innerhalb der Ebene selbst, denn
  der Standard rät von \emph{must} ab, während RFC~2119 \emph{MUST} und
  \emph{SHALL} gleichsetzt und fünf Stufen statt einer kennt. Er ist
  auflösbar, sobald eine der beiden Konventionen gewählt wird.
  \emph{Gegenargument (spräche für Note~2):} In der Ausprägung nach dem
  Standard ist die Ebene mit den Ebenen~(b) und~(d) widerspruchsfrei.
  \emph{Hinweis:} Diese Wahl ist inzwischen zugunsten der
  Schlüsselwörter des Standards gefallen
  (Abschnitt~\ref{T-subsec:rollenzuweisung}); die Note bleibt dennoch
  bei~1, weil sie die Ebene als Ganzes bewertet und die Bewertung auf
  dem Stand bleiben soll, auf dem sie in die Entscheidung eingegangen
  ist.

  \item[(d) Qualitäts-, Sprach- und Attributregeln, Note~2.]
  \emph{Beleg:} \cite[Abschn.~5.2.4 NOTE~2 und
  5.2.5--5.2.7]{iso29148}.
  \emph{Begründung:} Die Regeln beurteilen einen fertigen Satz und
  schreiben keine Satzform vor; sie kollidieren deshalb mit keiner
  Satzform. Der Standard nennt User Stories sogar ausdrücklich als
  zulässige alternative Formulierung.
  \emph{Gegenargument (spräche für Note~1):} Vollständige Konformität
  mit dem Standard setzt voraus, dass die Abschnitte~5.2.4 bis 5.2.7
  gemeinsam eingehalten werden \cite[Abschn.~4.2]{iso29148}, und für die
  Inhaltsvorgaben schließt er ein Anpassen bei vollständiger Konformität
  ausdrücklich aus \cite[Anhang~C.1]{iso29148}; eine Verbindung, die nur
  einen Teil dieser Vorgaben übernimmt, gibt den Anspruch auf
  Konformität auf.
\end{description}



\section{Tabellen zum Quality-User-Story-Rahmenwerk}
\label{sec:anhang-qus}

\begin{table}[htbp]
  \centering
  \caption{Die dreizehn Kriterien des Quality-User-Story-Rahmenwerks
           nach \cite[Tab.~1]{lucassen2016qus}}
  \label{tab:qus-kriterien}
  \begin{tabular}{@{}>{\raggedright\arraybackslash}p{3.1cm}%
                    >{\raggedright\arraybackslash}p{7.5cm}%
                    >{\raggedright\arraybackslash}p{1.9cm}@{}}
    \toprule
    \textbf{Kriterium} & \textbf{Bedeutung} & \textbf{Bezug} \\
    \midrule
    \multicolumn{3}{@{}l}{\textit{Syntaktisch}} \\
    Well-formed & enthält mindestens eine Rolle und ein Means & einzeln \\
    Atomic & beschreibt genau ein Merkmal & einzeln \\
    Minimal & enthält nichts außer Rolle, Means und Ends & einzeln \\
    \midrule
    \multicolumn{3}{@{}l}{\textit{Semantisch}} \\
    Conceptually sound & das Means benennt ein Merkmal, das End eine
      Begründung & einzeln \\
    Problem-oriented & benennt nur das Problem, nicht dessen Lösung &
      einzeln \\
    Unambiguous & vermeidet Begriffe, die mehrere Deutungen zulassen &
      einzeln \\
    Conflict-free & steht zu keiner anderen Story im Widerspruch &
      Menge \\
    \midrule
    \multicolumn{3}{@{}l}{\textit{Pragmatisch}} \\
    Full sentence & ist ein wohlgeformter vollständiger Satz & einzeln \\
    Estimatable & ist nicht so grob gefasst, dass Planung und
      Priorisierung erschwert werden & einzeln \\
    Unique & kommt kein zweites Mal vor & Menge \\
    Uniform & verwendet dieselbe Schablone wie die übrigen Stories &
      Menge \\
    Independent & ist in sich abgeschlossen, ohne inhärente
      Abhängigkeiten & Menge \\
    Complete & die Umsetzung der Menge ergibt eine funktionsvollständige
      Anwendung & Menge \\
    \bottomrule
  \end{tabular}
\end{table}

\begin{table}[htbp]
  \centering
  \caption{Trefferquote und Genauigkeit der selbsttätigen
           Qualitätsprüfung von User Stories, ermittelt an
           1023~Stories aus achtzehn Datensätzen
           \cite[Tab.~4 und 5]{lucassen2016qus}}
  \label{tab:aqusa}
  \begin{tabular}{@{}>{\raggedright\arraybackslash}p{4.5cm}rr@{}}
    \toprule
    \textbf{Kriterium} & \textbf{Trefferquote} & \textbf{Genauigkeit} \\
    \midrule
    Well-formed & 95,7\,\% & 42,3\,\% \\
    Atomic      & 82,9\,\% & 51,2\,\% \\
    Minimal     & 95,5\,\% & 69,5\,\% \\
    Uniform     & 95,4\,\% & 87,1\,\% \\
    Unique      & 100\,\%  & 100\,\% \\
    \midrule
    Gesamt (gewichtetes Mittel) & 93,8\,\% & 72,2\,\% \\
    \bottomrule
  \end{tabular}
\end{table}

\bibliographystyle{plain}
\begin{btSect}{thesis}
\section*{Literaturverzeichnis}
\btPrintCited
\end{btSect}
\begin{btSect}{online}
\section*{Online-Quellen}
\btPrintCited
\end{btSect}
\end{document}
